\subsection*{Addendum: the projective image over
\texorpdfstring{$\mathbb Z/4$}{Z/4}}

\begin{statement}{C.12}{Proved}
{The remaining finite projective image has order twenty-four}
Let $R=\mathbb Z/4$. Then

$$
\operatorname{im}\bigl(A_p(R)\longrightarrow\operatorname{PGL}_2(R)\bigr)
 =
 \operatorname{im}\bigl(\operatorname{SL}_2(R)
                  \longrightarrow\operatorname{PGL}_2(R)\bigr).
$$

This image has order $24$. Its complement is exactly the $24$
projective classes of nonsquare determinant listed in the
$R=\mathbb Z/4$ row of the replacement computation. In particular
$[\operatorname{diag}(3,1)]$ does not lift through $A_p(R)$.
This settles the finite image question left open in C.11.
The other source-section and integral-point questions stated there
are not settled by this addendum.
\end{statement}

\begin{proof}
We use the big fppf sheaf, which can be evaluated on arbitrary
$R$-algebras. In particular a localization not finitely presented
over $R$ is a legitimate object on which to pull back a section;
it is not thereby asserted to be an fppf cover.

\emph{The source on flat objects.}
Write $t=u(1)$ in
$\Gamma_p=\operatorname{PSL}_2(\mathbb Z)$.
Since $4=0$ in every additive $R$-group, the pushout defining $A_p$
factors the constant group through

$$
H=\Gamma_p/\langle\!\langle t^4\rangle\!\rangle
  =\operatorname{SL}_2(R)/\{\pm I\}.
$$

The last equality, including order $24$, is C.7. Its cusp
$C=\langle t\rangle$ has order four. The universal property of
pushouts therefore identifies the base-changed source with
$\underline H*_{\underline C}\mathbb G_a$.

Restricting to flat $R$-objects commutes with sheafification:
every fppf cover of a flat $R$-object still consists of flat
$R$-objects, so the same local constructions compute the restriction.
On this restricted site $\underline C\to\mathbb G_a$ is a
monomorphism. Indeed every nonzero flat algebra over the Artinian
local ring $R$ is faithfully flat. The same applies to any nonzero
clopen factor, so distinct locally constant residues modulo four
cannot acquire the same additive value. This assertion is not used
on the nonflat $R$-object $\mathbb F_2$.

\emph{Local preservation of ordinary normal forms.}
For every nonzero local flat $R$-algebra $S$, put
$W(S)=H*_C(S,+)$. The natural map $W(S)\to A_p(S)$ is injective.
To prove this, suppose an ordinary reduced word becomes the identity
in the sheaf. Equality in sheafification is witnessed locally by
equality in the presheaf, after refining a cover if necessary.
Choose a covering point above the closed point of
$\operatorname{Spec}S$ and localize its covering algebra there.
The resulting local map $S\to S'$ is flat and hence faithfully flat
\citep[Lemma~10.39.17, Tag~00HR]{C_Stacks_00HR}.
In particular it is injective. Since both rings are local, their
constant $H$ and $C$ sections are exactly $H$ and $C$.
An additive syllable outside $C$ cannot become an element of $C$
in $S'$, and the constant syllables do not change. The ordinary
amalgam normal-form theorem gives a contradiction.

\emph{A fixed reduced word on an Artinian local refinement.}
Take $\alpha\in A_p(R)$, and let
$g\in\operatorname{PGL}_2(R)$ be its matrix image.
Choose a nonempty affine fppf cover
$\operatorname{Spec}B\to\operatorname{Spec}R$ on which $\alpha$
is represented by a finite word. We may choose $B$ finitely
presented over $R$, so it is Noetherian. Let $P$ be a minimal prime
of $B$ and put $B_0=B_P$. This is a nonzero Noetherian local ring
of dimension zero, hence Artinian. It is flat and faithfully flat
over $R$. The map
$\operatorname{Spec}B_0\to\operatorname{Spec}R$
is an affine fpqc cover; it need not be fppf.

Reduce the pulled-back word in $W(B_0)$. If it is the identity,
then $g=1$ by fpqc separatedness of the representable target.
Otherwise write it as $x_1\cdots x_m$, with fixed alternating
factor types. For $m\geq2$ every syllable lies outside $C$.
For $m=1$ the subsequent argument uses the trivial insertion group.

Set $D=B_0\otimes_R B_0$. For every prime $q$ of $D$, the ring
$D_q$ is nonzero and flat over either copy of $B_0$.
Because $B_0$ is Artinian local, each map $B_0\to D_q$ is local,
and is therefore faithfully flat by Tag~00HR. Equivalently one
may use flatness over an Artinian local ring
\citep[Lemma~10.101.2, Tag~051G]{C_Stacks_051G}.
Consequently both pullback words retain the same reduced length
and factor types over $D_q$. Their products are equal in $A_p(D_q)$,
since both come from $\alpha$. The local injectivity just proved
makes them equal as ordinary reduced words.

\emph{Unique finite insertions and their descent.}
Equality of reduced words of these fixed types is equivalent to
unique elements $c_1,\ldots,c_{m-1}\in C$ satisfying

$$
y_1=x_1c_1,\qquad
 y_i=c_{i-1}^{-1}x_i c_i\quad(1<i<m),\qquad
 y_m=c_{m-1}^{-1}x_m.
$$

For clarity, uniqueness and existence follow by cancellation:
$x_1^{-1}y_1$ must lie in $C$, or the full inverse-product word
would remain reduced; this determines $c_1$, and induction gives
the rest. At each $D_q$ the inserted elements are ordinary
constant residues. There are finitely many additive equalities
in the displayed formulas. Each equality holding in $D_q$ holds
on some principal open neighbourhood of $q$. The constant-group
equalities hold there too. Thus the insertion tuples exist
Zariski locally on $\operatorname{Spec}D$.

Let $G=\underline C^{\,m-1}$ and let $X$ be the product of the
$m$ indicated factors, each either $\underline H$ or
$\mathbb G_a$. This is an affine $R$-scheme. There is a right action

$$
(x_1,\ldots,x_m)\cdot(c_1,\ldots,c_{m-1})
 =
 (x_1c_1,c_1^{-1}x_2c_2,\ldots,c_{m-1}^{-1}x_m).
$$

It is free on flat $R$-objects: equality of first coordinates
forces $c_1=1$, and induction forces every $c_i=1$, using the
monomorphism $C\to\mathbb G_a$ only on flat objects.
Therefore the local insertion tuples agree on intersections and
glue to $c\in G(D)$. On the triple overlap, freeness gives
$c_{13}=c_{12}c_{23}$.

No globally decidable word-length stratum has been introduced:
the fixed word was obtained over $B_0$, its reducedness was
checked in the local rings of overlaps, and the finite equalities
were then spread to open neighbourhoods.

The cocycle descends the trivial $G$-torsor over $B_0$ to a finite
\'etale $G$-torsor $T\to\operatorname{Spec}R$.
Effective descent of affine morphisms is
\citet[Lemma~35.37.1, Tag~0245]{C_Stacks_0245};
finite locally free algebras descend by faithfully flat module
descent, and \'etaleness descends by
\citet[Lemma~35.23.31, Tag~02VN]{C_Stacks_02VN}.
The tuple $(x_1,\ldots,x_m)$ gives a compatible equivariant map
from that trivial torsor to $X$, so it descends to a morphism
$f:T\to X$. Equivalently, $f$ is a tuple of ordinary word
coordinates over the finite \'etale cover $T$.

\emph{Only the matrix image is descended back.}
Let $\theta:X\to\operatorname{PGL}_2$ multiply the matrix images
of the coordinates. Cusp insertions cancel, so $\theta$ is
$G$-invariant. Consequently $\theta\circ f$ agrees with $g_T$
after the faithfully flat pullback
$B_0\otimes_R\mathcal O(T)$.
Since scheme-representable functors are fpqc sheaves
\citep[Lemma~35.13.7, Tag~023Q]{C_Stacks_023Q},
the two maps agree already over $T$.

Each $H$-coordinate has a determinant-one matrix representative
over $R$; choosing one for each of the finitely many elements of
$H$ also gives representatives for locally constant
$H$-coordinates over $T$. The additive coordinates have their
elementary determinant-one representatives. Their product is
therefore an element of $\operatorname{SL}_2(T)$ representing
$g_T$.

This conclusion uses fpqc descent for the affine coordinates and
the matrix target. It does not claim fpqc descent for $A_p$, nor
that $\alpha_T$ itself has been identified with the descended
word inside $A_p(T)$.

\emph{The determinant obstruction survives every such cover.}
The element $3$ is not a square in any nonzero finite \'etale
$R$-algebra $T$. Indeed $T/2T$ is reduced. If $z^2=3$, then
$(\overline z-1)^2=0$, so $\overline z=1$.
Hence $z=1+2b$ for some $b\in T$, and $z^2=1$, a contradiction
to $3\ne1$ in the faithfully flat algebra $T$.

If $g=[\operatorname{diag}(3,1)]$ had arisen from $A_p(R)$,
the preceding construction would give a determinant-one matrix
representing it over $T$. Equality of two matrix representatives
in $\operatorname{PGL}_2(T)$ means that their ratio is a global
scalar unit. Taking determinants would make $3$ a square in $T$,
which is impossible.

Finally C.10 already supplies the whole projective special-linear
subgroup, of order $24$ and index two. Any additional image
element would yield the excluded class after multiplication by
an element of that subgroup. Thus the image is exactly the
stated subgroup.
\end{proof}

\paragraph{Effect on the earlier residual.}
In C.11, the uncertainty about the order of the projective image
over $\mathbb Z/4$ and the lifting of
$[\operatorname{diag}(3,1)]$ is superseded by C.12:
that class does not lift, and the image has order $24$.
The computed list of $24$ missing word-image classes is now also
the complete list of missing sheaf-image classes.
The countably infinite source point group in C.8 is unchanged.
Complete source-group descriptions and the stated integral-point
questions remain outside this addendum.
